%% sections/related_work.tex
\section{Related Literature}
\label{sec:related}

The simulator sits at the intersection of four strands of
literature: the experimental-bubble tradition, the heterogeneous
agent-based finance tradition, the expected-utility decision-theoretic
tradition, and the emerging work on AI agents as
economic~subjects.

\paragraph{Experimental bubbles.} The taxonomy of the market
in this paper is directly inherited from \citet{smith1988bubbles},
who introduced the declining-fundamental, finite-horizon asset and
documented the characteristic bubble-and-crash trajectory in a
continuous double auction. \citet{dufwenberg2005bubbles} show that
the Smith--Suchanek--Williams bubble is fragile to \emph{experience}
rather than to instructions, incentives, or trading institutions:
once two-thirds of the subjects have played the game before,
transaction prices track the fundamental staircase from the first
period onward. The testbed in this paper inherits the DLM~2005
market substrate---asset life, dividend distribution, fundamental
staircase, four-round session loop---without modification, and
reuses the endogenous experience counter $k_i$ to ramp Plan~I's
belief-blend weight across the session, mirroring the paper's
procedural-experience channel at the level of the agent's belief
update rather than the agent's decision rule.

\paragraph{Heterogeneous agent-based finance.} Three classical
rules supply the pinned background of the simulator: a
\emph{Fundamentalist}~(F) that mean-reverts toward $\FV_t$ in the
Smith--Suchanek--Williams lineage; a \emph{Trend~follower}~(T) that
implements the positive-feedback strategy of
\citet{delong1990positive} by extrapolating a six-tick slope; and
a \emph{Random}~(R) zero-intelligence trader after
\citet{gode1993allocative} that posts uniform bids on
$[0.7,1.1]\,\FV_t$ and uniform asks on $[0.9,1.3]\,\FV_t$.
Together these three rules provide a rational anchor, a momentum
amplifier, and a liquidity floor; the utility agent layers an
expected-utility optimiser on top and is the class whose
belief-update plan the present paper varies.

\paragraph{Expected utility and AI-agent traders.}
\citet{lopezlira2025llm} introduces a utility-maximising AI-agent
trader whose prior is a perturbation of~$\FV_t$, whose action space is
$\{\text{hold},\text{cross bid},\text{cross ask},\text{passive bid},\text{passive ask}\}$,
and whose choice rule picks the action with the highest expected
utility,
$
\EU(a)=p_{\text{fill}}\cdot U(w_1)+(1-p_{\text{fill}})\cdot U(w_0).
$
Agents carry one of three risk-preference functionals
$\{\Uloving(w),\Uneutral(w),\Uaverse(w)\}=
\{(w/w_0)^2,\,w/w_0,\,\sqrt{w/w_0}\}$ and exchange private valuations
through a lightweight messaging bus. A pairwise trust score
$\tau_{r\to s}\in[0,1]$ is updated by exponential moving average
whenever a claimed valuation is observed against the period's
volume-weighted average trade price. The present paper inherits
this framework exactly as described---the prior, the action set,
the EU score, the risk functionals, the trust EMA, and the
messaging bus are all reimplemented without modification---and
then varies only the period-boundary belief update by comparing
the algorithmic blend of Plan~I against the language-model calls
of Plan~II and~III.

\paragraph{Language models in economic decision tasks.}
A small but rapidly growing literature studies whether language
models reproduce well-documented behavioural regularities when
embedded in economic decision tasks, and whether their outputs are
sensitive to prompt structure. The contrast between Plan~II and
Plan~III is a specific instance of this broader question: Plan~II
hands the model the closed-form expression of the risk-utility
functional $U_L, U_N, U_A$; Plan~III hands the model only the risk
label. Both prompts route through the same clamp, the same
consume-and-cache pipeline, and the same expected-utility scoring
rule downstream, so any observed difference in the posterior
distribution or in downstream market statistics is attributable to
the prompt structure alone. This narrow A/B design makes
prompt-sensitivity measurable on the same scale as the canonical
bubble statistics.

\paragraph{Market-quality instruments.} The testbed reports four
standard bubble statistics---Haessel's~$R^2$
\citep{haessel1978measuring}, normalised absolute deviation,
amplitude, and turnover, following SSW and
\citet{stockl2010bubble}---alongside allocative efficiency,
\begin{equation*}
\AllocEff\;=\;\frac{\sum_i V_i^{\text{post}}\cdot\text{inventory}_i}
{\max_i V_i^{\text{post}}\cdot Q},
\end{equation*}
which quantifies the fraction of theoretically-maximal utility
realised at the end of the session given the final posterior
valuations. All five instruments are reported for every plan on
every run, so cross-plan comparisons can be made without any
extra code.
